% Cover letter for Environment and Security (Sage).
% Submitted alongside the anonymized manuscript and the separate title page.
\documentclass[11pt]{article}
\usepackage[margin=1in]{geometry}
\usepackage{parskip}
\pagestyle{empty}
\begin{document}

\noindent Dear Editors,

We are pleased to submit ``Critical Minerals, Governance and Armed Conflict.
Evidence from Africa'' for consideration at \textit{Environment and Security}.

The paper asks whether the institutions governing mineral extraction change the
violence that follows discovery. Using geocoded records of critical mineral
discoveries across Africa from 1997 to 2018, combined with UCDP event data at
seven spatial resolutions, we estimate staggered difference-in-differences models
for three minerals chosen because each sits under a different governance regime.
Tantalum falls under international due diligence frameworks such as Dodd-Frank
Section 1502 and ITSCI, platinum group metals under South Africa's Mineral and
Petroleum Resources Development Act, and chromium under neither.

Holding region and period fixed while varying the governance regime lets us treat
institutional design as the explanatory variable rather than as a control, which
is not how most of this literature is built. We have three main findings. Discovery
reduces the frequency of conflict under every regime, including the absence of
one. Only national governance also reduces lethality. International due diligence
removes organized armed groups without any matching improvement for civilians,
and where no framework applies, conflict events fall while battle deaths rise,
with state forces the only armed actor present.

An actor-level decomposition using the UCDP Actor Dataset shows why. Governance
changes which actors are present around a deposit rather than reducing violence
as such. What appears to separate the regimes is not how much governance exists
but whether state control over extraction is itself constrained.

We believe that this paper is a good fit for this journal because it speaks to conflict
financing through natural resources and the political ecology of extraction. It
also relates to work the journal has published on the governance of critical
minerals, including Deberdt and Le Billon on the limits of techno-legal fixes,
whose argument we test. That literature is largely
institutional and legal and rarely asks whether these instruments change violence
on the ground, which is the gap we address.

The finding has a practical implication that evaluations of supply chain due
diligence that use the presence or activity of armed groups as their outcome will
record success in settings where civilian security has not improved. 

The manuscript is approximately 7,900 words excluding references, tables and
figures. It is not under consideration elsewhere, and all data and replication
materials that may be shared will be made available on publication.

\vspace{1em}
\noindent Yours sincerely,

\vspace{1em}
\noindent Elisa D'Amico, University College Dublin \\
Päivi Lujala, University of Oulu

\end{document}
